\subsection{Introduction}
In the previous chapter we have seen an overview of the LQR problem and how to solve it using dynamic programming. In that formulation, we assumed that there were no constraints on the state and control input. However, in many real-world applications, constraints on the state and input must be explicitly taken into account.\\
In the presence of constraints, the problem cannot, in general, be solved in closed form via backward induction as in the unconstrained LQR case, since the value function is no longer quadratic.The resulting optimization problem becomes:

\begin{theorembox}[Constrained LQ finite horizon problem]
\begin{align*}
\min_{u_0, \ldots, u_{T-1},x_0,\dots,x_T} &\sum_{t=0}^{T-1} \left( x_t^\top Q x_t + u_t^\top R u_t \right) + x_T^\top S x_T \\
\text{s.t.} \quad & x_{t+1} = A x_t + B u_t, \quad t=0, \ldots, T-1 \\
& x_0 = x_0 \\
& \textcolor{red}{x_t \in \mathcal{X}_t}\\
& \textcolor{red}{u_t \in \mathcal{U}_t}
\end{align*}
Usually, the constraint sets are convex, e.g., polytopes.
\end{theorembox}

One could solve this optimization problem in a \textbf{one-shot} fashion (i.e., open-loop strategy), obtaining an optimal input sequence $u_0, \ldots, u_{T-1}$ for the given initial state $x_0$, instead of a state-feedback policy as in the unconstrained case (where $u_t = \Gamma_t x_t$).\\ However, if disturbances or model mismatch cause the system to deviate from the predicted trajectory, the previously computed control sequence is no longer optimal. As a result, the system will follow a different trajectory than expected (see Figures \ref{fig:mpc_01} and \ref{fig:mpc_02}).\\
To address this issue, the optimal control problem must be re-solved using the updated state information. Since we assume a Markovian system, the current state contains all the information needed to compute an optimal control sequence for the future. This implies that the optimization problem must be solved repeatedly, at each time step, within one sampling interval.

\begin{figure}[!ht]
\centering
\begin{tikzpicture}[>=Stealth, line cap=round, line join=round, scale=1.0]

    % Axes
    \draw[->, thick] (0,0) -- (11.5,0) node[below right] {$t$};
    \draw[->, thick] (0,0) -- (0,3.2) node[left] {$u$};

    % Horizon marker
    \draw[dashed, thick] (1.6,-0.05) -- (1.6,2.7);

    % T mark
    \draw[thick] (10.0,-0.15) -- (10.0,0.25);
    \node[below] at (10.05,-0.1) {$T$};

    % Legend
    \node[draw=gray, circle, minimum size=5pt, inner sep=0pt] at (6.0,2.55) {};
    \node[anchor=west] at (6.4,2.55) {previously planned input};

    \node[draw=red, circle, minimum size=5pt, inner sep=0pt] at (6.0,2.05) {};
    \node[anchor=west] at (6.4,2.05) {planned input};

    \node[fill=red, draw=red, circle, minimum size=5pt, inner sep=0pt] at (6.0,1.55) {};
    \node[anchor=west] at (6.4,1.55) {applied input};

    % Previously planned input (gray open circles)
    \foreach \x/\y in {
        1.6/2.55,
        2.05/1.25,
        2.45/0.45,
        3.15/0.00,
        3.55/0.00,
        3.95/0.00,
        4.35/0.00
    }{
        \draw[gray] (\x,\y) circle (0.07);
    }

    % Planned input (red open circles)
    \foreach \x/\y in {
        1.6/2.20,
        2.1/1.95,
        2.9/1.60,
        3.4/0.60,
        3.9/0.95,
        4.35/0.95,
        4.85/0.30,
        5.3/0.05,
        5.75/0.03,
        6.2/0.03,
        6.65/0.03,
        7.1/0.03,
        7.55/0.03,
        8.0/0.03,
        8.45/0.03,
        8.9/0.03,
        9.35/0.03,
        10.0/0.03
    }{
        \draw[red] (\x,\y) circle (0.07);
    }

    % Applied input (red filled circles)
    \foreach \x/\y in {
        0.0/1.35,
        0.5/1.95,
        1.0/2.50,
        1.35/1.95
    }{
        \fill[red] (\x,\y) circle (0.07);
    }

\end{tikzpicture}
\caption{Receding-horizon update of the input sequence: previously planned inputs are shifted, a new optimal input sequence is computed based on the current state.}
\label{fig:mpc_01}
\end{figure}

\begin{figure}[!ht]
\centering
\begin{tikzpicture}[>=Stealth, line cap=round, line join=round, scale=1.0]
    % Axes
    \draw[->, thick] (0,-0.6) -- (0,2.5) node[left] {$x$};
    \draw[->, thick] (-0.5,0) -- (11.3,0) node[below right] {$t$};

    % "now" dashed line
    \draw[dashed, thick] (1.8,-0.05) -- (1.8,2.4);
    \node[below] at (1.8,-0.12) {\large now};

    % Horizon marker T
    \draw[thick] (9.0,-0.18) -- (9.0,0.18);
    \node[below] at (9.35,-0.08) {\Large $T$};

    % Legend
    \draw[gray] (5.6,1.95) circle (0.07);
    \node[anchor=west] at (6.0,1.95) {previously predicted state};

    \draw[blue!80!black] (5.6,1.45) circle (0.07);
    \node[anchor=west] at (6.0,1.45) {predicted state};

    \fill[blue!80!black] (5.6,0.95) circle (0.07);
    \node[anchor=west] at (6.0,0.95) {measured state};

    % Previously predicted state (gray open circles)
    \foreach \x/\y in {
        0.90/1.5,
        1.35/1,
        1.80/0.55,
        2.25/0.28,
        2.70/0.55,
        3.15/0.00,
        3.60/0.00,
        4.05/0.00
    }{
        \draw[gray] (\x,\y) circle (0.07);
    }

    % Predicted state (blue open circles)
    \foreach \x/\y in {
        2.25/1.05,
        2.70/1.05,
        3.15/0.80,
        3.60/0.52,
        4.05/0.25,
        4.50/0.02,
        4.95/0.02,
        5.40/0.02,
        5.85/0.02,
        6.30/0.02,
        6.75/0.02,
        7.20/0.02,
        7.65/0.02,
        8.10/0.02,
        8.55/0.02,
        9.00/0.02
    }{
        \draw[blue!80!black] (\x,\y) circle (0.07);
    }

    % Measured state (blue filled circles)
    \foreach \x/\y in {
        0.00/1.95,
        0.45/1.95,
        0.90/1.70,
        1.35/1.42,
        1.80/1.15
    }{
        \fill[blue!80!black] (\x,\y) circle (0.07);
    }

\end{tikzpicture}
\caption{Receding-horizon update of the state trajectory: previously predicted states are corrected using the measured state at the current time, and a new optimal state trajectory is predicted over the horizon.}
\label{fig:mpc_02}
\end{figure}

The idea that emerges from this discussion is to solve the optimization problem in a \textbf{receding horizon} fashion. At each time step, we solve a finite-horizon optimal control problem of length $T$, and apply only the first control input of the optimal sequence (see Figure \ref{fig:mpc_03}).\\
At the next time step, the new state is measured, and the optimization problem is solved again with the updated initial condition. If the system is time-invariant, the optimization problem remains the same at each time step, except for the initial state. The corresponding predicted state evolution over the horizon is illustrated in Figure \ref{fig:mpc_04}.\\

\begin{figure}[!ht]
\centering
\begin{tikzpicture}[>=Stealth, line cap=round, line join=round, scale=1.0]
    % Axes
    \draw[->, thick] (-2.3,0) -- (8.2,0) node[below right] {$t$};
    \draw[->, thick] (0,-0.6) -- (0,2.6) node[left] {$u$};

    % Horizon marker T
    \draw[thick] (3.2,-0.2) -- (3.2,0.2);
    \node at (1.5,0.35) {\Large $T$};

    % Labels past / future
    \node at (-1.2,-0.45) {\large past};
    \node at (1.2,-0.45) {\large future};

    % Applied inputs (past - red filled)
    \foreach \x/\y in {
        -2.0/0.55,
        -1.6/0.55,
        -1.2/0.80,
        -0.8/0.55,
        -0.4/1.05
    }{
        \fill[red] (\x,\y) circle (0.08);
    }

    % Current input at time 0 (highlighted)
    \fill[green!70!black] (0,1.10) circle (0.08);

    % Planned inputs (future - red open)
    \foreach \x/\y in {
        0.8/1.55,
        1.2/2.20,
        1.6/1.55,
        2.2/2.20,
        2.6/1.00
    }{
        \draw[red] (\x,\y) circle (0.08);
    }

\end{tikzpicture}
\caption{Finite-horizon input sequence in receding horizon control: at the current time, an optimal sequence of future inputs over a horizon of length $T$ is computed, while past inputs are fixed.}
\label{fig:mpc_03}
\end{figure}

\begin{figure}[!ht]
\centering
\begin{tikzpicture}[>=Stealth, line cap=round, line join=round, scale=1.0]

    % Axes
    \draw[->, thick] (-2.3,0) -- (8.7,0) node[below right] {$t$};
    \draw[->, thick] (0,-0.6) -- (0,2.6) node[left] {$x$};

    % Labels past / future
    \node at (-1.2,-0.45) {\large past};
    \node at (1.2,-0.45) {\large future};

    % Horizon marker T
    \draw[thick] (3.2,-0.2) -- (3.2,0.2);
    \node at (1.5,0.35) {\Large $T$};

    % Measured states (past - filled blue)
    \foreach \x/\y in {
        -2.0/1.4,
        -1.5/2.0,
        -1.0/1.4,
        -0.5/1.4,
        0.0/2.0
    }{
        \fill[blue!80!black] (\x,\y) circle (0.08);
    }

    % Predicted states (future - open blue)
    \foreach \x/\y in {
        0.5/2.0,
        1.0/1.4,
        1.5/0.9,
        2.0/0.5,
        2.5/0.0
    }{
        \draw[blue!80!black] (\x,\y) circle (0.08);
    }

\end{tikzpicture}
\caption{Finite-horizon state prediction in receding horizon control: starting from the current state, future states are predicted over a horizon of length $T$ using the planned input sequence.}
\label{fig:mpc_04}
\end{figure}

Thus the following scheme defines the receding horizon control strategy, also known as \textbf{Model Predictive Control (MPC)}:
\begin{enumerate}
\item At time $t$, measure the current state $x_t$.
\item Solve the finite-horizon optimal control problem with initial state $x_t$ to obtain the optimal input sequence $u_t, u_{t+1}, \ldots, u_{t+T-1}$.
\item Apply the first control input $u_t$ to the system.
\item Increment time $t \leftarrow t+1$ and repeat from step 1.
\end{enumerate}

The resulting \textbf{open-loop optimal control problem} corresponds to a \textit{parametric
optimization problem} parametrized in the \textcolor{red}{initial state x}.
\begin{theorembox}[Parametric optimization problem ]
$u^*(x)$ is determined by\footnote{Notation: $k$ is the index for the optimization problem, while $t$ is the index for the time steps of the system.}
\begin{align*}
    \min_{\mathbf{u},\mathbf{x}} &\sum_{k=0}^{T-1} \left( x_k^\top Q x_k + u_k^\top R u_k \right) + x_T^\top S x_T\\
    \text{s.t.} \quad & x_{k+1} = A x_k + B u_k\\
    & \textcolor{red}{x_0 = x} \\
    & x_k \in \mathcal{X}_k \\
    & u_k \in \mathcal{U}_k
\end{align*}
\end{theorembox}

The parametric nature of the problem immediately raises the question of whether the resulting control strategy is \emph{feedforward} or \emph{feedback}. Although the optimization problem computes an open-loop input sequence, the receding horizon implementation introduces feedback implicitly. Indeed, at each time step the optimization problem is solved using the current state $x_t$, and only the first input of the optimal sequence is applied. Therefore, the implemented control law can be written as
\[
u_t = u^*(x_t),
\]
which defines a \textbf{state-feedback control law}. One can have a visual intuition of this feedback structure by looking at the block diagram in Figure \ref{fig:mpc_05}.
\begin{figure}[!ht]
\centering
\begin{tikzpicture}[
    >=Stealth,
    block/.style={draw, thick, minimum width=3cm, minimum height=1.8cm},
    node distance=4cm
]

\node[block] (controller) {$u^*(\cdot)$};
\node[block, right=of controller] (plant) {plant};

\draw[->, thick] (controller.east) -- node[above] {$u_t$} (plant.west);
\draw[->, thick] (plant.east) -- ++(1.5,0) node[above] {$x_t$};

\coordinate (out) at ($(plant.east)+(1.5,0)$);
\coordinate (low) at ($(out)+(0,-2)$);
\coordinate (leftlow) at ([xshift=-1.2cm]controller.west |- low);
\coordinate (leftmid) at ([xshift=-1.2cm]controller.west);

\draw[->, thick]
    (out) -- (low) -- (leftlow) -- (leftmid) -- (controller.west);

\end{tikzpicture}

\caption{Block diagram of the MPC control scheme}
\label{fig:mpc_05}
\end{figure}

We can establish some properties about this optimization problem:
\begin{itemize}
    \item Under \textbf{continuity} assumptions on the cost and the dynamics, the minimum cost is a constinuous function of the initial state $x$.
    \item Under \textbf{compactness} assumptions on the constraints, the parametric optimization problem admits a solution for all $x$ in the feasible set.
    \item The map $x \mapsto u^*(x)$ is not necessarily continuos for a general cost function
\end{itemize}
In the case of linear dynamics, quadratic cost, and convex constraints, the optimization problem is a convex quadratic program (QP) parametrized in the initial state $x$. In this case, the map $x \mapsto u^*(x)$ is piecewise affine and continuous. However, in general, the control law induced by MPC can be nonlinear and non-smooth.
\begin{theorembox}[MPC control law]
As a consequence, MPC can be interpreted as a \textbf{memory-less, nonlinear, time-invariant feedback control law}. The nonlinearity arises from the presence of constraints, even when the system dynamics are linear and the cost is quadratic.\\
\end{theorembox}

A natural question is whether MPC can achieve \emph{perfect tracking}, i.e., zero steady-state error. In general, this is not guaranteed without additional design choices (e.g., terminal ingredients or disturbance modeling), since the controller repeatedly solves a finite-horizon problem and does not explicitly enforce asymptotic properties beyond the prediction horizon.
To better understand the relation between MPC and the finite-horizon optimal control problem, consider the following example:
\paragraph{Example}
\begin{align*}
\min_{\mathbf{x},\mathbf{u}} \quad & \sum_{t=0}^{T-1} a^t |u_t| + |x_T|^2, \quad |a| < 1 \\
\text{s.t.} \quad & x_{t+1} = x_t + u_t.
\end{align*}

Two different questions can be asked:
\begin{enumerate}
    \item What is the optimal open-loop solution $(\mathbf{x}^*, \mathbf{u}^*)$?
    \item What is the MPC feedback law $x \mapsto u^*(x)$ obtained by solving the problem in a receding horizon fashion?
\end{enumerate}
To solve this we observe that the cost is not quadratic and that the cost of applying inouts is discounted over time. Furthermore the state is penalized only at the end of the horizon.\\
Thus the optimal solution is to simply wait for the last time step and then apply $u_{T} = -x_{T}$, which yields a cost of $|x_0|^2$.\\
In contrast, the MPC feedback law is given by the control law $x \mapsto u^*(x)=0$.\\
A key observation is that the optimal trajectory associated with the finite-horizon problem is, in general, \textbf{never realized in closed loop}, even in the absence of disturbances or model mismatch. This is because at each time step the optimization problem is re-solved with a shifted horizon, leading to a different optimal sequence than the one computed previously.
$\blacksquare$

\medskip

Consider now the unconstrained linear-quadratic problem:
\begin{align*}
\min_{\mathbf{x},\mathbf{u}} \quad & \sum_{t=0}^{T-1} \left( x_t^\top Q x_t + u_t^\top R u_t \right) + x_T^\top S x_T, \quad Q,S \succeq 0,\; R \succ 0 \\
\text{s.t.} \quad & x_{t+1} = A x_t + B u_t.
\end{align*}
\textbf{Finite-time LQR} yields a time-varying linear feedback law
\begin{align*}
    u_t &= \Gamma_t x_t\\
    \Gamma_t &= -(R + B^\top P_{t+1} B)^{-1} B^\top P_{t+1} A
\end{align*}
where the matrices $P_t$ are computed backward in time. The control gain explicitly depends on time through $P_{t+1}$.\\
In contrast, \textbf{MPC} repeatedly solves the same finite-horizon problem at each time step. As a result, the applied control law is
\[
u_t = \Gamma_0 x_t,
\]
where $\Gamma_0$ is the gain associated with the first step of the horizon. Hence, MPC induces a \textbf{linear time-invariant} feedback law.\\
More explicitly, at time $t=0$ both approaches yield the same control input:
\[
u_0^{\text{MPC}} = -(R + B^\top P_1 B)^{-1} B^\top P_1 A x_0 = u_0^{\text{LQR}}.
\]
However, at the next time step a discrepancy appears. Finite-time LQR applies
\[
u_1^{\text{LQR}} = \textcolor{blue}{-(R + B^\top P_2 B)^{-1} B^\top P_2 A} x_1,
\]
while MPC recomputes the problem and applies again
\[
u_1^{\text{MPC}} = {-(R + B^\top \textcolor{red}{P_1} B)^{-1} B^\top \textcolor{red}{P_1} A} x_1.
\]
Therefore, even in the unconstrained case, MPC does not reproduce the finite-horizon optimal control sequence beyond the first step.

\subsection{Implementation of MPC}
The implementation of MPC requires solving an optimization problem at each time step. We show two main approaches to implement MPC:
\begin{enumerate}
    \item Compute the \textbf{function} $u^*(\cdot)$ \textbf{offline}, and then evaluate it online at each time step.
    \begin{itemize}
        \item It is extremely complex to solve with exception of very few cases
        \item Abundant offline computational resources are required 
        \item Large memory is required to store the function $u^*(\cdot)$
    \end{itemize}
    \item Compute the \textbf{value} of the function $U^*(\cdot)$ \textbf{online} at each time step by solving the optimization problem.
    \begin{itemize}
        \item It is ogten a tractable optimization problem.
        \item It has limited time to perform computation online.
    \end{itemize}
\end{enumerate}

\subsubsection{Explicit MPC}

Explicit MPC corresponds to the first implementation strategy: instead of solving the optimization problem online at each sampling instant, we try to compute the feedback law$ x \mapsto u^*(x) $
offline.\\
We have already seen one important case. In the \textbf{unconstrained linear MPC} case, the optimizer is simply the first element of the finite-horizon LQR solution. Therefore, the applied input can be written as
\[
u_t = u^*(x_t)=\Gamma_0 x_t,
\]
so the control law is linear.\\

The natural question is what happens when \emph{constraints} are introduced. Consider the standard constrained linear MPC problem

\begin{theorembox}[Linear MPC with linear constraints]
\begin{align*}
\min_{\mathbf{u},\mathbf{x}} \quad &
\sum_{k=0}^{T-1}\left(x_k^\top Q x_k + u_k^\top R u_k\right) + x_T^\top S x_T \\
\text{s.t.}\quad
& x_{k+1}=Ax_k+Bu_k, \qquad k=0,\dots,T-1, \\
& x_0=x, \\
& x_k \in \mathcal{X}, \qquad \mathcal{X}=\{x \in \mathbb{R}^n : Cx \le d\}, \\
& u_k \in \mathcal{U}, \qquad \mathcal{U}=\{u \in \mathbb{R}^m : Eu \le f\},
\end{align*}
with $ Q,S \succeq 0, \qquad R \succ 0 $.\\
This is the most common MPC formulation: linear dynamics, quadratic cost, and polyhedral state/input constraints.
\end{theorembox}
This problem is a \textbf{convex quadratic program} parametrized by the initial state $x$. In general, the optimizer is no longer globally linear. However, due to the quadratic cost and linear constraints, it still has a very specific structure: the explicit MPC law turns out to be \textbf{piecewise affine}. The idea is best understood through a very simple example.

\paragraph{Example}
Consider a scalar integrator with horizon $T=1$:
\begin{align*}
\min_{u_0,x_1} \quad & u_0^2 + x_1^2 \\
\text{s.t.}\quad & x_1 = x_0 + u_0, \\
& x_1 \le 1.
\end{align*}

This example is useful because:
\begin{itemize}
    \item it is a quadratic problem with a linear constraint,
    \item it is parametrized by the current state $x_0$,
    \item the equality constraint allows us to eliminate the variable $x_1$,
    \item the problem is convex, hence the KKT conditions are necessary and sufficient.
\end{itemize}
Using the model equation $x_1=x_0+u_0$, we can eliminate $x_1$ and rewrite the problem as
\[
\min_{u_0} f(u_0)
\qquad \text{s.t.} \qquad g(u_0)\le 0,
\]
where
\begin{align*}
f(u_0) &= u_0^2 + (x_0+u_0)^2, \\
g(u_0) &= x_0 + u_0 - 1.
\end{align*}
The derivatives are
\begin{align*}
\nabla f(u_0) &= 4u_0 + 2x_0, \\
\nabla g(u_0) &= 1.
\end{align*}

\begin{theorembox}[KKT conditions for one inequality constraint]
For a problem of the form
\[
\min_x f(x)
\qquad \text{s.t.}\qquad g(x)\le 0,
\]
the KKT conditions are
\begin{align*}
\nabla f(x)+\mu \nabla g(x) &= 0, \qquad \mu \ge 0,\\
g(x) &\le 0,\\
\mu\, g(x) &= 0.
\end{align*}
The last relation is the \emph{complementary slackness} condition.
\end{theorembox}

Applying the KKT conditions to the problem gives
\begin{align*}
4u_0 + 2x_0 + \mu &= 0, \qquad \mu \ge 0, \\
x_0 + u_0 - 1 &\le 0, \\
\mu(x_0 + u_0 - 1) &= 0.
\end{align*}

The complementary slackness condition implies that two cases are possible.

\paragraph{Case 1: inactive constraint.}
Assume first that the inequality is inactive, so
\[
\mu = 0.
\]
Then stationarity becomes
\[
4u_0 + 2x_0 = 0,
\]
hence
\[
u_0 = -\frac{x_0}{2}.
\]
To check whether this candidate is feasible, we substitute it into the constraint:
\[
x_0 + u_0 - 1 = x_0 - \frac{x_0}{2} - 1 = \frac{x_0}{2} - 1.
\]
Therefore feasibility requires
\[
\frac{x_0}{2} - 1 \le 0
\qquad \Longleftrightarrow \qquad
x_0 \le 2.
\]
So, for all initial states satisfying $x_0 \le 2$, the optimal input is
\[
u_0^* = -\frac{x_0}{2}.
\]

\paragraph{Case 2: active constraint.}
Assume now that the inequality is active, so
\[
x_0 + u_0 - 1 = 0.
\]
Then
\[
u_0 = 1 - x_0.
\]
Substituting this into the stationarity condition,
\[
4u_0 + 2x_0 + \mu = 0,
\]
we obtain
\[
\mu = -4u_0 - 2x_0 = -4(1-x_0)-2x_0 = 2x_0 - 4.
\]
Since $\mu \ge 0$, this case is valid only if
\[
2x_0 - 4 \ge 0
\qquad \Longleftrightarrow \qquad
x_0 \ge 2.
\]
Hence, for all $x_0 \ge 2$, the optimal input is
\[
u_0^* = 1 - x_0.
\]

Combining the two cases, the explicit solution is
\begin{theorembox}[Explicit solution of the toy MPC problem]
\[
u_0^*(x_0)=
\begin{cases}
-\dfrac{x_0}{2}, & x_0 \le 2, \\[1ex]
1-x_0, & x_0 \ge 2.
\end{cases}
\]
\end{theorembox}

This example shows the main structural property of explicit MPC:
\begin{itemize}
    \item the control law is \textbf{continuous},
    \item the control law is \textbf{piecewise affine},
    \item each affine expression is valid in a region where the same set of constraints is active.
\end{itemize}

The reason is the following. Once the active constraints are known, the KKT conditions reduce to a linear system in the decision variables and the multipliers. Solving that linear system gives an optimizer that depends \emph{affinely} on the parameter $x$. When the active set changes, a different affine expression is obtained. Therefore, the state space is partitioned into regions, and on each region the controller is affine.\\

Although the explicit solution is conceptually appealing, it becomes difficult to use in large problems. In practice:
\begin{itemize}
    \item there may be tens or hundreds of constraints,
    \item the number of regions may become very large,
    \item the offline computation required to determine all regions can be very long,
    \item even online, one must determine in which region the current state lies before evaluating the affine law.
\end{itemize}
As a consequence, explicit MPC is mainly attractive when the system dimension is small and the sampling time is so short that solving a quadratic program online would be too expensive. For larger systems, the most common approach is instead to solve the MPC optimization problem online at each sampling instant.

\subsection{Stability}

A useful way to think about MPC is as a \emph{feedback design tool}. Once the prediction horizon $T$ and the weights $Q,R,S$ are chosen, the receding-horizon implementation induces a control law of the form
\[
u_t = u^*(x_t).
\]
Even though the optimization problem computes an open-loop input sequence, only the first input is applied and the problem is solved again at the next time step. Therefore, the implemented controller is a \textbf{static, time-invariant, nonlinear feedback law}. The nonlinearity does not necessarily come from the dynamics; it already appears because of the repeated constrained optimization.\\
From this perspective, the main question is no longer only how to solve the optimization problem, but whether the resulting closed-loop system is stable. There are two conceptually different ways to approach this question. The first is an \emph{analysis} point of view: one tries to compute the control law $u^*(\cdot)$ explicitly and then study the stability of the closed-loop dynamics. In general, this is difficult, because the MPC law can be piecewise affine or more generally nonlinear and non-smooth. The second is a \emph{design} point of view: instead of computing the feedback law explicitly, one looks for conditions on the design parameters $T,Q,R,S$ that guarantee closed-loop stability directly. This is the more useful viewpoint in practice.\\

To study stability, the standard tool is Lyapunov analysis. Consider a discrete-time system
\[
x_{t+1}=f(x_t),
\]
and suppose that the origin $x=0$ is an equilibrium, that is, $f(0)=0$.

\begin{theorembox}[Stable equilibrium]
The equilibrium $x=0$ is said to be \textbf{stable} if
\[
\forall \varepsilon > 0, \ \exists \delta > 0 \text{ such that } \norm{x_0}<\delta
\ \Longrightarrow\ 
\norm{x_t}<\varepsilon \quad \forall t\geq 0.
\]
In words, if the initial condition is sufficiently close to the equilibrium, then the entire trajectory remains close to it for all future times.
\end{theorembox}

Stability alone only guarantees that trajectories do not move far away from the equilibrium. A stronger notion is asymptotic stability, which additionally requires convergence to the equilibrium.

\begin{theorembox}[Asymptotically stable equilibrium]
The equilibrium $x=0$ is \textbf{asymptotically stable} if it is stable and
\[
\lim_{t\to\infty}\norm{x_t}=0.
\]
Hence, trajectories starting sufficiently close to the origin remain close to it and eventually converge to it.
\end{theorembox}

The classical criterion to prove asymptotic stability is based on a Lyapunov function.

\begin{theorembox}[Lyapunov theorem for discrete-time systems]
Assume there exists a real-valued function $W:\mathbb{R}^n\to\mathbb{R}$ such that
\begin{align*}
W(0) &= 0,\\
W(x) &> 0 \qquad \forall x\neq 0,\\
W(f(x)) &< W(x) \qquad \forall x\neq 0.
\end{align*}
Then the equilibrium $x=0$ is asymptotically stable.
\end{theorembox}

The idea is simple: $W$ plays the role of an energy-like quantity. It is zero only at the equilibrium, positive everywhere else, and strictly decreases along the closed-loop trajectories. Therefore, the system is forced to move toward the origin.\\

This Lyapunov viewpoint explains why stability is relatively transparent for \emph{infinite-horizon} MPC. Suppose that, at time $t$, the controller computes an input sequence that minimizes the infinite-horizon cost
\[
V_t^\infty(x,\mathbf{u}) = \sum_{k=t}^{\infty} g_k(x_k,u_k),
\]
subject to the system dynamics. If the minimum is attained, it is natural to define the optimal value function
\[
W(x) = \min_{\mathbf{u}} V_t^\infty(x,\mathbf{u}).
\]
Now comes the key observation. By Bellman's principle of optimality, if an input sequence is optimal from time $t$, then its tail is optimal from time $t+1$ onward. In other words, after applying the first optimal input $u^*(x_t)$ and reaching the next state $x_{t+1}$, the remaining part of the optimal sequence is still optimal for the optimization problem starting at $t+1$. As a consequence, along the closed-loop trajectory one obtains
\[
W(x_{t+1}) = W(x_t) - g_t(x_t,u^*(x_t)).
\]
Therefore, if the stage cost satisfies suitable positivity properties, for example if
\[
g_t(x_t,u_t) > 0 \qquad \text{for all } x_t\neq 0,
\]
then
\[
W(x_{t+1}) < W(x_t) \qquad \text{for all } x_t\neq 0.
\]
This means that the optimal value function is a natural Lyapunov candidate for the closed-loop system. Under reasonable assumptions, it follows that the origin is asymptotically stable.\\

This argument is very elegant, but it also hides some subtleties. First, an infinite-horizon MPC problem is infinite-dimensional, so from a computational viewpoint it is not directly implementable as stated. Second, one is implicitly assuming that the global minimum exists and is attained. Third, the stage cost must satisfy suitable properties so that the value function is positive definite and decreases strictly along the closed loop. These requirements are analogous in spirit to the assumptions encountered in infinite-horizon LQR.\\

The situation changes significantly for \emph{finite-horizon} MPC. In that case, the same Lyapunov argument does not apply automatically. At time $t$, the controller minimizes a cost over the interval $[t,t+T]$. At time $t+1$, however, the optimization problem is not simply the tail of the previous one: the first stage cost has disappeared, but a new term has entered at the end of the horizon. Therefore, the controller at time $t+1$ is minimizing a \emph{different} objective. As a result, the optimizer changes, and the finite-horizon value function is not automatically decreasing along the closed-loop trajectory.

This is the fundamental reason why closed-loop stability is immediate for infinite-horizon optimal control, but not for finite-horizon receding-horizon MPC. For finite-horizon MPC, stability must be enforced through additional design choices. This will be the starting point for the next part of the discussion.